% =====================================================================
% BMED365: Computational Modeling
% Beamer Presentation - Topic List Z01-Z13
% =====================================================================
\documentclass[aspectratio=169, 10pt]{beamer}

% =====================================================================
% PACKAGES
% =====================================================================
\usepackage[utf8]{inputenc}
\usepackage[T1]{fontenc}
\usepackage[english]{babel}
\usepackage{graphicx}
\usepackage{tikz}
\usetikzlibrary{shapes.geometric, arrows, positioning, calc}
\usepackage{booktabs}
\usepackage{amsmath}
\usepackage{fontawesome5}
\usepackage{hyperref}
\hypersetup{colorlinks=true, linkcolor=blue, urlcolor=blue}

% =====================================================================
% THEME AND COLORS
% =====================================================================
\usetheme{Madrid}
\usecolortheme{spruce}

% Custom colors
\definecolor{modelingteal}{RGB}{0, 128, 128}
\setbeamercolor{structure}{fg=modelingteal}

% =====================================================================
% TITLE INFO
% =====================================================================
\title{Computational Modeling}
\subtitle{BMED365: Topic List Z01--Z13}
\author{BMED365}
\date{Spring 2026}
\institute{Department of Biomedicine, University of Bergen}

% =====================================================================
% DOCUMENT
% =====================================================================
\begin{document}

% Title page
\begin{frame}
    \titlepage
\end{frame}

% Table of Contents
\begin{frame}{Overview}
    \tableofcontents
\end{frame}

% =====================================================================
% SECTION: Introduction to Computational Modeling
% =====================================================================
\section{Fundamentals of Computational Modeling}

\begin{frame}{Z01: Explain the role of mathematical models in biomedicine}
    \textbf{Why Mathematical Models in Biomedicine?}
    
    \vspace{0.3cm}
    \begin{block}{Definition}
        A \textbf{mathematical model} is a description of a system using mathematical concepts and language -- enabling simulation, prediction, and understanding of biological processes.
    \end{block}
    
    \vspace{0.3cm}
    \begin{columns}[T]
        \begin{column}{0.48\textwidth}
            \textbf{Key Roles:}
            \begin{itemize}
                \item \textbf{Understanding:} Reveal mechanisms
                \item \textbf{Prediction:} Forecast outcomes
                \item \textbf{Optimization:} Design interventions
                \item \textbf{Integration:} Combine data types
                \item \textbf{Hypothesis testing:} Virtual experiments
            \end{itemize}
        \end{column}
        \begin{column}{0.48\textwidth}
            \textbf{Applications:}
            \begin{itemize}
                \item Drug pharmacokinetics
                \item Disease progression
                \item Cardiac electrophysiology
                \item Tumor growth dynamics
                \item Neural activity
                \item Epidemiology
            \end{itemize}
        \end{column}
    \end{columns}
    
    \vspace{0.3cm}
    \begin{alertblock}{The Digital Twin Vision}
        Patient-specific models could enable personalized treatment simulation before actual intervention.
    \end{alertblock}
\end{frame}

\begin{frame}{Z02: Distinguish between mechanistic and phenomenological models}
    \textbf{Two Fundamental Modeling Approaches}
    
    \vspace{0.3cm}
    \begin{columns}[T]
        \begin{column}{0.48\textwidth}
            \textbf{Mechanistic Models:}
            \begin{itemize}
                \item Based on physical/biological principles
                \item Parameters have biological meaning
                \item Derived from first principles
                \item More complex, harder to fit
                \item Better extrapolation
            \end{itemize}
            
            \vspace{0.2cm}
            \textbf{Example:}\\
            Hodgkin-Huxley model -- ion channels, conductances, Nernst potentials
        \end{column}
        \begin{column}{0.48\textwidth}
            \textbf{Phenomenological Models:}
            \begin{itemize}
                \item Describe observed behavior
                \item Parameters may lack meaning
                \item Fitted to data empirically
                \item Simpler, easier to use
                \item Limited extrapolation
            \end{itemize}
            
            \vspace{0.2cm}
            \textbf{Example:}\\
            Logistic growth -- $\frac{dN}{dt} = rN(1-\frac{N}{K})$
        \end{column}
    \end{columns}
    
    \vspace{0.3cm}
    \begin{block}{Practical Choice}
        Use \textbf{mechanistic} models when you need to understand \textit{why}; use \textbf{phenomenological} models when you need to predict \textit{what}.
    \end{block}
\end{frame}

\begin{frame}{Z03: Describe ordinary differential equations (ODEs) in biological systems}
    \textbf{ODEs: The Language of Dynamic Systems}
    
    \vspace{0.3cm}
    \begin{block}{What is an ODE?}
        An equation relating a function to its derivatives:\\[0.2cm]
        $\displaystyle \frac{dy}{dt} = f(y, t, \theta)$\\[0.2cm]
        where $y$ is the state, $t$ is time, and $\theta$ are parameters.
    \end{block}
    
    \vspace{0.3cm}
    \begin{columns}[T]
        \begin{column}{0.48\textwidth}
            \textbf{Common in biology:}
            \begin{itemize}
                \item Population dynamics
                \item Chemical kinetics
                \item Drug metabolism
                \item Gene regulation
                \item Neural firing
            \end{itemize}
        \end{column}
        \begin{column}{0.48\textwidth}
            \textbf{Solving ODEs:}
            \begin{itemize}
                \item Analytical (rare in biology)
                \item Numerical integration
                \item Euler, Runge-Kutta methods
                \item Python: \texttt{scipy.integrate}
            \end{itemize}
        \end{column}
    \end{columns}
    
    \vspace{0.3cm}
    \begin{alertblock}{Key Insight}
        Many biological processes involve \textbf{feedback loops} and \textbf{nonlinearities} -- ODEs capture these dynamics naturally.
    \end{alertblock}
\end{frame}

% =====================================================================
% SECTION: Electrophysiology
% =====================================================================
\section{Neural and Cardiac Electrophysiology}

\begin{frame}{Z04: Explain the Hodgkin-Huxley model of action potentials}
    \textbf{The Nobel Prize-Winning Model (1963)}
    
    \vspace{0.2cm}
    \begin{block}{Core Idea}
        The action potential arises from voltage-dependent opening and closing of ion channels (Na$^+$, K$^+$) in the cell membrane.
    \end{block}
    
    \vspace{0.2cm}
    \begin{columns}[T]
        \begin{column}{0.48\textwidth}
            \textbf{The HH Equations:}
            \begin{align*}
                C_m \frac{dV}{dt} &= I - g_{Na}m^3h(V-E_{Na})\\
                &\quad - g_K n^4(V-E_K) - g_L(V-E_L)
            \end{align*}
            
            \textbf{Components:}
            \begin{itemize}
                \item $V$: membrane potential
                \item $m, h, n$: gating variables
                \item $g$: conductances
                \item $E$: reversal potentials
            \end{itemize}
        \end{column}
        \begin{column}{0.48\textwidth}
            \textbf{Action Potential Phases:}
            \begin{enumerate}
                \item \textbf{Resting:} $V \approx -70$ mV
                \item \textbf{Depolarization:} Na$^+$ influx
                \item \textbf{Repolarization:} K$^+$ efflux
                \item \textbf{Hyperpolarization:} Overshoot
                \item \textbf{Recovery:} Return to rest
            \end{enumerate}
            
            \vspace{0.2cm}
            \begin{alertblock}{\footnotesize Lab 5}
                \footnotesize Notebook 01: Simulate HH model and explore parameter effects.
            \end{alertblock}
        \end{column}
    \end{columns}
\end{frame}

\begin{frame}{Z07--Z08: Cardiovascular flow and waveforms}
    \textbf{Modeling Blood Flow and Cardiac Signals}
    
    \vspace{0.3cm}
    \begin{columns}[T]
        \begin{column}{0.48\textwidth}
            \textbf{Z07: Cardiovascular Flow Models}
            
            \vspace{0.2cm}
            \textbf{Windkessel Model:}
            \begin{itemize}
                \item Arteries as elastic reservoir
                \item Resistance + Compliance
                \item Predicts pressure waveforms
            \end{itemize}
            
            \vspace{0.2cm}
            \textbf{1D Blood Flow:}
            \begin{itemize}
                \item Navier-Stokes simplification
                \item Wave propagation in arteries
                \item Pulse wave velocity
            \end{itemize}
        \end{column}
        \begin{column}{0.48\textwidth}
            \textbf{Z08: ECG and Blood Pressure}
            
            \vspace{0.2cm}
            \textbf{ECG Modeling:}
            \begin{itemize}
                \item Cardiac dipole theory
                \item P, QRS, T waves
                \item Lead configurations
            \end{itemize}
            
            \vspace{0.2cm}
            \textbf{Pressure Waveforms:}
            \begin{itemize}
                \item Systolic/diastolic peaks
                \item Dicrotic notch
                \item Arterial stiffness markers
            \end{itemize}
        \end{column}
    \end{columns}
    
    \vspace{0.3cm}
    \begin{block}{Lab 5 Notebook 03}
        Simulate cardiovascular models with ECG and blood pressure waveforms.
    \end{block}
\end{frame}

% =====================================================================
% SECTION: Tumor and Cell Biology
% =====================================================================
\section{Tumor Growth and Cell Biology}

\begin{frame}{Z05--Z06: Tumor growth and angiogenesis models}
    \textbf{Mathematical Oncology}
    
    \vspace{0.3cm}
    \begin{columns}[T]
        \begin{column}{0.48\textwidth}
            \textbf{Z05: Tumor Growth Models}
            
            \vspace{0.2cm}
            \textbf{Classic Models:}
            \begin{itemize}
                \item \textbf{Exponential:} $\frac{dN}{dt} = \lambda N$
                \item \textbf{Logistic:} $\frac{dN}{dt} = \lambda N(1-\frac{N}{K})$
                \item \textbf{Gompertz:} $\frac{dN}{dt} = \lambda N \ln(\frac{K}{N})$
            \end{itemize}
            
            \vspace{0.2cm}
            \textbf{Parameters:}
            \begin{itemize}
                \item $\lambda$: growth rate
                \item $K$: carrying capacity
                \item $N$: tumor size/cell count
            \end{itemize}
        \end{column}
        \begin{column}{0.48\textwidth}
            \textbf{Z06: Angiogenesis Models}
            
            \vspace{0.2cm}
            \textbf{Why it matters:}
            \begin{itemize}
                \item Tumors need blood supply
                \item Anti-angiogenic therapy
                \item Metastasis pathway
            \end{itemize}
            
            \vspace{0.2cm}
            \textbf{Modeling approaches:}
            \begin{itemize}
                \item VEGF diffusion
                \item Endothelial cell migration
                \item Vessel network formation
                \item Nutrient delivery
            \end{itemize}
        \end{column}
    \end{columns}
    
    \vspace{0.3cm}
    \begin{alertblock}{Clinical Relevance}
        Tumor growth models inform treatment scheduling, dose optimization, and predict response to therapy.
    \end{alertblock}
\end{frame}

\begin{frame}{Z09--Z10: Muscle force and cell signaling}
    \textbf{Biomechanics and Systems Biology}
    
    \vspace{0.3cm}
    \begin{columns}[T]
        \begin{column}{0.48\textwidth}
            \textbf{Z09: Muscle Force Models}
            
            \vspace{0.2cm}
            \textbf{Hill Model:}
            \begin{itemize}
                \item Force-velocity relationship
                \item Contractile + elastic elements
                \item Activation dynamics
            \end{itemize}
            
            \vspace{0.2cm}
            \textbf{Applications:}
            \begin{itemize}
                \item Movement biomechanics
                \item Prosthetics design
                \item Sports performance
                \item Cardiac muscle function
            \end{itemize}
        \end{column}
        \begin{column}{0.48\textwidth}
            \textbf{Z10: Cell Signaling Networks}
            
            \vspace{0.2cm}
            \textbf{ODE-based modeling:}
            \begin{itemize}
                \item Receptor-ligand binding
                \item Enzyme kinetics (Michaelis-Menten)
                \item Phosphorylation cascades
                \item Feedback loops
            \end{itemize}
            
            \vspace{0.2cm}
            \textbf{Gene Regulatory Networks:}
            \begin{itemize}
                \item Transcription factor dynamics
                \item Oscillations (circadian)
                \item Bistability (cell fate)
            \end{itemize}
        \end{column}
    \end{columns}
    
    \vspace{0.3cm}
    \begin{block}{Lab 5 Notebooks}
        04: Muscle force models \quad 05: Cell signaling and gene regulation
    \end{block}
\end{frame}

% =====================================================================
% SECTION: AI-Assisted Modeling
% =====================================================================
\section{AI-Assisted Modeling}

\begin{frame}{Z11: Describe how to use LLMs to assist model development}
    \textbf{LLMs as Modeling Assistants}
    
    \vspace{0.3cm}
    \begin{columns}[T]
        \begin{column}{0.48\textwidth}
            \textbf{How LLMs can help:}
            \begin{itemize}
                \item Generate model equations
                \item Write simulation code
                \item Explain model behavior
                \item Debug numerical issues
                \item Literature review
                \item Parameter estimation ideas
            \end{itemize}
            
            \vspace{0.2cm}
            \textbf{Example prompt:}\\
            \footnotesize\textit{``Write Python code to simulate the Hodgkin-Huxley model using scipy.integrate.solve\_ivp''}
        \end{column}
        \begin{column}{0.48\textwidth}
            \textbf{Best practices:}
            \begin{itemize}
                \item Verify generated equations
                \item Check units and dimensions
                \item Test with known solutions
                \item Iterate on prompts
                \item Combine with domain expertise
            \end{itemize}
            
            \vspace{0.2cm}
            \begin{alertblock}{Lab 5 Notebook 06}
                Use LLM to generate a kidney filtration model from scratch -- demonstrating AI-assisted modeling workflow.
            \end{alertblock}
        \end{column}
    \end{columns}
\end{frame}

\begin{frame}{Z12--Z13: Systems biology and model validation}
    \textbf{Advanced Modeling Concepts}
    
    \vspace{0.3cm}
    \begin{columns}[T]
        \begin{column}{0.48\textwidth}
            \textbf{Z12: Systems Biology Approaches}
            
            \vspace{0.2cm}
            \textbf{Multi-scale modeling:}
            \begin{itemize}
                \item Molecular $\rightarrow$ Cellular
                \item Cellular $\rightarrow$ Tissue
                \item Tissue $\rightarrow$ Organ
                \item Organ $\rightarrow$ Organism
            \end{itemize}
            
            \vspace{0.2cm}
            \textbf{Integration challenges:}
            \begin{itemize}
                \item Different time scales
                \item Parameter transfer
                \item Computational cost
            \end{itemize}
        \end{column}
        \begin{column}{0.48\textwidth}
            \textbf{Z13: Model Validation}
            
            \vspace{0.2cm}
            \textbf{Validation approaches:}
            \begin{itemize}
                \item Compare with experimental data
                \item Sensitivity analysis
                \item Uncertainty quantification
                \item Cross-validation
            \end{itemize}
            
            \vspace{0.2cm}
            \textbf{Key questions:}
            \begin{itemize}
                \item Does the model fit the data?
                \item Are predictions robust?
                \item What are the limitations?
            \end{itemize}
        \end{column}
    \end{columns}
    
    \vspace{0.3cm}
    \begin{block}{Resources}
        \textbf{BioModels:} Database of curated models -- \url{https://www.ebi.ac.uk/biomodels/}\\
        \textbf{CellML:} Model exchange format -- \url{https://www.cellml.org/}
    \end{block}
\end{frame}

% =====================================================================
% SUMMARY
% =====================================================================
\section*{Summary}

\begin{frame}{Summary: Computational Modeling}
    \textbf{Key Learning Points:}
    \begin{itemize}
        \item \textbf{Z01--Z03:} Fundamentals -- mechanistic vs. phenomenological, ODEs
        \item \textbf{Z04, Z07--Z08:} Electrophysiology -- Hodgkin-Huxley, cardiovascular
        \item \textbf{Z05--Z06:} Oncology -- tumor growth, angiogenesis
        \item \textbf{Z09--Z10:} Biomechanics -- muscle, cell signaling
        \item \textbf{Z11--Z13:} Advanced -- LLM-assisted, systems biology, validation
    \end{itemize}
    
    \vspace{0.3cm}
    \begin{block}{Lab 5 Notebooks}
        \footnotesize
        \begin{tabular}{ll}
            00-test-llm & Local LLM setup (DeepSeek-R1) \\
            01-action-potentials & Hodgkin-Huxley model \\
            02-tumor-growth & Tumor dynamics \\
            03-cardiovascular-flow & Blood flow + ECG \\
            04-muscle-force & Force generation \\
            05-cell-signaling & Gene networks \\
            06-kidney-filtration & LLM-assisted modeling \\
        \end{tabular}
    \end{block}
    
    \vspace{0.2cm}
    \begin{alertblock}{Individual Project}
        Consider modeling projects: cardiac, neural, tumor, pharmacokinetics...
    \end{alertblock}
\end{frame}

\end{document}

